Los capítulos anteriores dejaron dos deudas pendientes, ambas señaladas en su momento pero nunca saldadas. La primera: la Sección \ref{sec:ecualizacion canal zfe} planteó el filtro ecualizador de cero forzado y llegó hasta decir que ``da $2N+1$ ecuaciones para determinar las $2N+1$ ganancias $c_n$'', pero se detuvo ahí -- sin escribir el sistema, sin resolverlo, sin un ejemplo numérico. La segunda: los diagramas de bloques de los Capítulos \ref{chapter:deteccion ruido} y \ref{chapter:modulacion pasabanda} (Figuras \ref{fig:deteccion ruido} y \ref{fig:deteccion canal ruido}) vienen arrastrando una caja llamada ``Sincronismo'', conectada con una flecha punteada, que nunca se abre -- se supuso siempre, como dice explícitamente el texto que acompaña a la Figura \ref{fig:deteccion ruido}, que ``conocemos el momento exacto donde muestrear la onda recibida'', sin decir cómo se logra eso en la práctica.

Este capítulo cierra ambos huecos. Lo hace apoyándose en una pieza que antes no teníamos: el $\hat H(f)$ físico que el Capítulo \ref{chapter:propagacion guiada} le dio al canal guiado. La ecualización existe porque $\hat H(f)$ no es plano; el sincronismo existe porque el transmisor y el receptor no comparten reloj ni oscilador. Ninguno de los dos problemas es nuevo conceptualmente -- ya sabíamos que existían -- pero hasta ahora los tratamos como cajas negras.

\section{Ecualización de canal}\label{sec:ecualizacion}

\subsection*{Resolución explícita del filtro de cero forzado}

Recordemos el planteo de la Sección \ref{sec:ecualizacion canal zfe} (Figura \ref{fig:filtro eq}): se recibe una onda con ISI, $x_R(t) = \sum_k a_k \tilde p(t-kT_s)$, y se la pasa por un filtro transversal de $2N+1$ ganancias $\{c_n\}_{n=-N}^N$, obteniendo una nueva onda PAM con pulso
\begin{equation*}
    \hat p(t) = \sum_{n=-N}^N c_n\, \tilde p(t-nT_s).
\end{equation*}
El criterio de \emph{cero forzado} (ZFE) pide $\hat p(kT_s)=0$ para $k=\pm 1,\ldots,\pm N$, junto con $\hat p(0)=1$ (para no perder la ganancia de la muestra útil). Escribamos esto explícitamente: evaluando $\hat p$ en $t=kT_s$,
\begin{equation*}
    \hat p(kT_s) = \sum_{n=-N}^N c_n\, \tilde p\big((k-n)T_s\big) = \sum_{n=-N}^N c_n\, \tilde p_{k-n},
\end{equation*}
donde abreviamos $\tilde p_m := \tilde p(mT_s)$ (las muestras del pulso recibido, el ``dato experimental'' que menciona la Sección \ref{sec:ecualizacion canal zfe}). La condición $\hat p(kT_s)=\delta_{k,0}$ para $k=-N,\ldots,N$ es entonces un sistema \emph{lineal} de $2N+1$ ecuaciones en las $2N+1$ incógnitas $c_n$:
\begin{equation}\label{eq:sistema zfe}
    \sum_{n=-N}^N \tilde p_{k-n}\, c_n = \delta_{k,0}, \qquad k=-N,\ldots,N.
\end{equation}
En forma matricial, $\mathbf P\,\mathbf c = \mathbf e_0$, donde $\mathbf c=(c_{-N},\ldots,c_N)^\top$, $\mathbf e_0$ es el vector con un $1$ en la posición central y ceros en el resto, y $\mathbf P$ es la matriz $(2N+1)\times(2N+1)$ de entradas $P_{k,n}=\tilde p_{k-n}$. Como $P_{k,n}$ depende solo de $k-n$, $\mathbf P$ es una \emph{matriz de Toeplitz} (constante en cada diagonal). La solución formal es $\mathbf c = \mathbf P^{-1}\mathbf e_0$, que existe siempre que $\mathbf P$ sea inversible -- lo es en la generalidad de los casos de interés práctico.

\begin{ejemplo}[ZFE de tres ganancias, $N=1$]\label{ejemplo:zfe n1}
Supongamos que, al medir el canal, obtenemos las muestras del pulso recibido (simétricas, un caso frecuente en la práctica): $\tilde p_0=1$, $\tilde p_{\pm 1}=0{,}2$, $\tilde p_{\pm 2}=0{,}05$ (y $\tilde p_m\approx 0$ para $|m|>2$). Con $N=1$, el sistema \eqref{eq:sistema zfe} tiene tres ecuaciones ($k=-1,0,1$) en las incógnitas $c_{-1},c_0,c_1$:
\begin{equation*}
    \begin{pmatrix} \tilde p_0 & \tilde p_{-1} & \tilde p_{-2} \\ \tilde p_1 & \tilde p_0 & \tilde p_{-1} \\ \tilde p_2 & \tilde p_1 & \tilde p_0 \end{pmatrix}
    \begin{pmatrix} c_{-1} \\ c_0 \\ c_1 \end{pmatrix}
    =
    \begin{pmatrix} 1 & 0{,}2 & 0{,}05 \\ 0{,}2 & 1 & 0{,}2 \\ 0{,}05 & 0{,}2 & 1 \end{pmatrix}
    \begin{pmatrix} c_{-1} \\ c_0 \\ c_1 \end{pmatrix}
    =
    \begin{pmatrix} 0 \\ 1 \\ 0 \end{pmatrix}.
\end{equation*}
Por la simetría de $\tilde p$, la solución también es simétrica, $c_{-1}=c_1$: de la primera ecuación, $c_{-1}(1{,}05) + 0{,}2\, c_0 = 0 \Rightarrow c_{-1}=-0{,}1905\, c_0$; sustituyendo en la segunda, $0{,}4\,c_{-1}+c_0=1 \Rightarrow 0{,}9238\, c_0 = 1$. Resolviendo:
\begin{equation*}
    c_0 \approx 1{,}0825, \qquad c_{-1}=c_1\approx -0{,}2062.
\end{equation*}
Verificando: $\hat p(0)=1$, $\hat p(\pm T_s)=0$ exactamente, como se pidió. Pero $\hat p(\pm 2T_s)\approx 0{,}0129$: \emph{no} es cero -- el ZFE con $N=1$ solo garantiza los ceros dentro del rango pedido, no más allá. Para eliminar también esa ISI residual haría falta $N=2$ (cinco ganancias) o más, al costo de un filtro más largo.
\end{ejemplo}

Como ya anticipaba la Sección \ref{sec:ecualizacion canal zfe}, ajustar el filtro requiere conocer $\tilde p_m$, lo que se obtiene experimentalmente (transmitiendo un pulso o secuencia de entrenamiento conocida y midiendo la respuesta). En la práctica rara vez se calcula $\mathbf c=\mathbf P^{-1}\mathbf e_0$ de una sola vez: el canal puede variar levemente, o no conocerse con precisión de antemano, y es más robusto \emph{adaptar} los $c_n$ iterativamente a partir del error observado (algoritmos como LMS, \emph{Least Mean Squares}). No desarrollamos el algoritmo adaptativo en este curso -- el canal guiado es, a diferencia de uno inalámbrico, esencialmente \emph{invariante en el tiempo} una vez instalado, por lo que un ecualizador fijo, calibrado una vez, alcanza en la gran mayoría de los casos.

\subsection*{Ecualización en el dominio de la frecuencia}

El filtro transversal de la sección anterior fuerza un número \emph{finito} de ceros en el dominio del tiempo -- el Ejemplo \ref{ejemplo:zfe n1} mostró que eso deja ISI residual fuera del rango forzado. La alternativa conceptualmente más simple es trabajar directamente en frecuencia: si $\hat P(f)$ es el espectro del pulso recibido $\tilde p(t)$ (relacionado con $\hat H(f)$ del canal, Ecuación \eqref{eq:H canal}, y con el pulso y filtro de transmisión y recepción), un ecualizador de respuesta $\hat H_e(f)$ produce, a la salida, un pulso de espectro $\hat P(f)\hat H_e(f)$. Eligiendo
\begin{equation*}
    \hat H_e(f) := \frac{1}{\hat P(f)},
\end{equation*}
el pulso equalizado tiene espectro idénticamente $1$, es decir $\hat p(t)=\delta(t)$: \textbf{cero ISI exacto}, no solo en un rango finito de instantes de muestreo sino en todo tiempo. Esto es, en esencia, el mismo problema que el filtro transversal de la sección anterior -- invertir el efecto del canal -- pero visto en el dominio dual: la ecualización en tiempo resuelve un sistema lineal finito (convolución discreta con $2N+1$ términos); la ecualización en frecuencia es, formalmente, una simple división, $\hat H_e(f)=1/\hat P(f)$, válida en todo $f$. Los filtros transversales de la sección anterior son, de hecho, la forma práctica de \emph{aproximar} este cociente ideal con un número finito de coeficientes.

\subsection*{El costo de invertir en banda ancha}

La simplicidad de $\hat H_e(f)=1/\hat P(f)$ esconde un problema serio, que ya el texto original de la Sección \ref{sec:ecualizacion canal zfe} anticipaba cualitativamente (``si algunas ganancias son muy altas, el ruido puede verse amplificado''): en cualquier frecuencia donde $|\hat P(f)|$ sea chico -- y por la forma de $\hat H(f)$ que dedujimos en el Capítulo \ref{chapter:propagacion guiada} (Figura \ref{fig:atenuacion comparada}), $|\hat P(f)|$ decrece sin cota al crecer $f$ -- el ecualizador necesita una ganancia $|\hat H_e(f)|=1/|\hat P(f)|$ enorme. El problema es que el ruido a la entrada del ecualizador (Capítulo \ref{chapter:deteccion ruido}) tiene energía en todas las frecuencias, y esa misma ganancia enorme amplifica el ruido en esa banda exactamente tanto como reconstruye la señal. En el límite, si $\hat P(f)\to 0$ en alguna banda (un cero real de $\hat H(f)$, no solo una atenuación fuerte), $\hat H_e(f)\to\infty$ ahí: el ZFE en frecuencia intenta reconstruir información que, en esa banda, ya se perdió irremediablemente en ruido.

Esto sugiere una idea distinta, que retomamos con toda su fuerza en el Capítulo \ref{chapter:ofdm}: en vez de invertir $\hat P(f)$ de una sola vez sobre todo el ancho de banda, \emph{partir} la banda en tramos angostos, cada uno lo suficientemente chico como para que $\hat P(f)$ sea aproximadamente \emph{plano} dentro de él. En cada tramo, invertir $\hat P(f)$ ya no requiere una ganancia extrema -- es, esencialmente, dividir por una constante (un único número complejo por tramo) en vez de por una función que varía mucho. La idea completa -- banda ancha partida en muchos subcanales angostos, cada uno ecualizado trivialmente -- es exactamente lo que hace OFDM, y es el motivo por el cual, adelantándolo aquí, ese capítulo empieza a tener sentido: \textbf{OFDM es, en el fondo, una forma práctica de resolver el problema de ecualización que acabamos de plantear}, evitando el costo de invertir en banda ancha.

\section{Sincronización}\label{sec:sincronizacion}

\subsection*{El problema del sincronismo}

Repasemos qué venimos suponiendo, sin decirlo, desde hace varios capítulos. El detector coherente de la Sección \ref{sec:deteccion coherente} multiplica la señal recibida por $2\cos(2\pi f_c t)$ y $-2\sin(2\pi f_c t)$ -- osciladores locales de frecuencia \emph{exactamente} $f_c$ y fase \emph{exactamente} alineada con el oscilador que generó la portadora en el transmisor. Y el detector de la Figura \ref{fig:deteccion ruido} muestrea la salida del filtro de recepción en el instante $kT_s$ -- suponiendo que \emph{conocemos} ese instante con precisión. Ninguna de las dos cosas es gratis: el transmisor y el receptor son dos dispositivos físicamente separados, cada uno con su propio oscilador de referencia (un cristal de cuarzo, típicamente) y su propio reloj de muestreo, sin ningún cable en común que los sincronice.

Aunque los osciladores de transmisor y receptor estén diseñados para la misma frecuencia nominal $f_c$, ningún cristal real oscila a una frecuencia exacta -- hay una tolerancia de fabricación (típicamente unas pocas partes por millón). El oscilador local del receptor tiene entonces una frecuencia $f_c+\Delta f$, con $\Delta f$ un \emph{offset de frecuencia de portadora} (\emph{Carrier Frequency Offset}, CFO) desconocido pero acotado, y una fase inicial arbitraria. Lo mismo ocurre, de forma independiente, con el reloj que define los instantes de muestreo $kT_s$: puede estar corrido en fase, o incluso avanzar a una velocidad ligeramente distinta a $f_s$. Sin corregir estos desajustes, ni la fase de la constelación ni el instante de muestreo son los que el resto del texto supone.

El problema del sincronismo tiene, entonces, tres capas, que conviene distinguir porque se resuelven con mecanismos distintos y en un orden lógico:
\begin{enumerate}
    \item \textbf{Sincronización de portadora}: estimar y corregir el CFO y la fase del oscilador local, para que la demodulación coherente (Sección \ref{sec:deteccion coherente}) entregue realmente $x_i(t)\approx\sum a_k p(t-kT_s)$ y no una versión rotada y batiendo en el tiempo.
    \item \textbf{Temporización de símbolo}: estimar el instante correcto de muestreo $kT_s$, para que la muestra tomada corresponda al máximo del pulso (o, más generalmente, al punto de ISI nula que garantiza el pulso de Nyquist, Sección \ref{sec:nyquist no ideal}).
    \item \textbf{Sincronización de trama}: una vez enganchados portadora y reloj de símbolo, saber \emph{dónde empieza el mensaje} dentro del flujo continuo de símbolos -- y, como veremos, resolver de paso una ambigüedad de fase que las dos capas anteriores dejan sin resolver.
\end{enumerate}
Las primeras dos capas comparten una misma herramienta matemática -- un \emph{lazo de enganche de fase} (\emph{Phase-Locked Loop}, PLL) -- aplicada a dos variables físicas distintas (fase de portadora, fase de reloj de símbolo). Conviene entonces derivar el lazo genérico una sola vez, y después particularizarlo dos veces.

\subsection*{El lazo de enganche de fase (PLL)}

Un PLL es un sistema realimentado que ajusta la fase de un oscilador local hasta que coincide con la de una referencia entrante. Tres bloques lo componen (Figura \ref{fig:pll generico}):
\begin{itemize}
    \item Un \emph{detector de fase}, que compara la fase $\theta_{in}(t)$ de la señal entrante con la fase $\theta_{vco}(t)$ del oscilador local, y produce una señal de error $e(t)$ proporcional (para error chico) a la diferencia: $e(t) \approx K_d\,\big(\theta_{in}(t)-\theta_{vco}(t)\big)$, con $K_d$ la ganancia del detector.
    \item Un \emph{filtro de lazo}, que suaviza $e(t)$ -- en su versión más simple, una ganancia pura, que es la que consideramos aquí (\emph{lazo de primer orden}).
    \item Un \emph{oscilador controlado por tensión} (VCO), cuya frecuencia instantánea es proporcional a la tensión de control $v(t)$ que recibe: $\displaystyle\frac{d\theta_{vco}}{dt} = K_{vco}\, v(t)$, con $K_{vco}$ la sensibilidad del VCO. Es decir, el VCO \emph{integra} su entrada para producir una fase.
\end{itemize}

\begin{figure}
    \begin{center}
    \begin{blockdiagram}
        \node [etiqueta] (inicio) {$\theta_{in}(t)$};
        \node [operator, right=1.3cm of inicio] (sum) {$+$};
        \node [block, right=2.6cm of sum, align=center] (fd) {Detector\\de fase};
        \node [block, right=3.4cm of fd, align=center] (vco) {VCO\\[2pt]\footnotesize$\frac{d\theta_{vco}}{dt}=K_{vco}v(t)$};
        \node [etiqueta, right=1.8cm of vco] (salida) {$\theta_{vco}(t)$};
        \draw [->] (inicio) -- (sum);
        \draw [->] (sum) -- node[midway,above] {\footnotesize $\theta_e(t)$} (fd);
        \draw [->] (fd) -- node[midway,above] {\footnotesize $e(t)\approx K_d\theta_e(t)$} (vco);
        \draw [->] (vco) -- (salida);
        \draw [->] (vco.south) -- ++(0,-1.2) -| node[near start, below, xshift=-3.2cm]{\footnotesize $-\theta_{vco}(t)$} (sum.south);
    \end{blockdiagram}
    \end{center}
    \caption{Lazo de enganche de fase (PLL) genérico, de primer orden: el detector de fase compara la referencia entrante con la realimentación del propio VCO, y la diferencia corrige la frecuencia de este último hasta anular el error.}\label{fig:pll generico}
\end{figure}

Definamos el error de fase $\theta_e(t) := \theta_{in}(t) - \theta_{vco}(t)$. Derivando y sustituyendo las dos relaciones anteriores ($v(t)=e(t)\approx K_d\theta_e(t)$ para un lazo de primer orden sin filtro adicional, y la ecuación del VCO):
\begin{equation*}
    \frac{d\theta_e}{dt} = \frac{d\theta_{in}}{dt} - \frac{d\theta_{vco}}{dt} = \frac{d\theta_{in}}{dt} - K_{vco}K_d\,\theta_e(t).
\end{equation*}
Llamando $K:=K_{vco}K_d$ a la \emph{ganancia de lazo}, obtenemos la ecuación diferencial que gobierna el enganche:
\begin{equation}\label{eq:pll edo}
    \frac{d\theta_e}{dt} + K\,\theta_e(t) = \frac{d\theta_{in}}{dt}.
\end{equation}
Es una ecuación diferencial lineal de primer orden, con la derivada de la fase entrante como entrada. Dos casos cubren lo que nos interesa:

\begin{itemize}
    \item \textbf{Salto de fase puro} (la fase entrante da un salto $\Delta\theta$ en $t=0$ y luego permanece a frecuencia nominal, $d\theta_{in}/dt=0$ para $t>0$): la ecuación homogénea $\dot\theta_e = -K\theta_e$ tiene solución $\theta_e(t)=\theta_e(0)\,\e^{-Kt}$. El error de fase decae exponencialmente a cero, con constante de tiempo $1/K$: el lazo \emph{engancha}. Cuanto mayor la ganancia $K$, más rápida la adquisición.
    \item \textbf{Offset de frecuencia constante} $\Delta f$ (la fase entrante es $\theta_{in}(t)=2\pi\Delta f\, t+\theta_0$, es decir $d\theta_{in}/dt = 2\pi\Delta f$ constante): en régimen permanente ($\dot\theta_e=0$), \eqref{eq:pll edo} da
    \begin{equation}\label{eq:pll error regimen}
        \theta_{e,\infty} = \frac{2\pi\Delta f}{K}.
    \end{equation}
    El lazo de primer orden sí logra seguir un offset de frecuencia constante -- el VCO efectivamente adopta esa frecuencia --, pero queda con un \emph{error de fase residual} proporcional a $\Delta f$ e inversamente proporcional a la ganancia $K$. Subir $K$ reduce este error, pero (aunque no lo desarrollamos aquí) también ensancha la banda de ruido que el lazo deja pasar: hay un compromiso entre velocidad de adquisición/error residual y inmunidad al ruido en la elección de $K$.
\end{itemize}

Con este comportamiento genérico -- adquisición exponencial, con constante de tiempo $1/K$, y error residual ante un offset de frecuencia -- vamos a reencontrarnos dos veces: primero enganchando la fase de la portadora, después enganchando la fase del reloj de muestreo. En ambos casos la ecuación \eqref{eq:pll edo} es la misma; lo que cambia es qué señal física hace de $\theta_{in}(t)$ y cómo se construye el detector de fase.

\subsection*{Sincronización de portadora: el lazo de Costas}

Apliquemos el PLL genérico al problema de la Sección \ref{sec:deteccion coherente}. Sea, para BPSK, $x_c(t) = \sum_k a_k\, p(t-kT_s)\cos(2\pi f_c t+\theta)$, con $a_k\in\{-A,A\}$ y $\theta$ la fase (desconocida) del oscilador transmisor. El receptor genera localmente $2\cos(2\pi f_c t+\theta_{vco})$ y $-2\sin(2\pi f_c t+\theta_{vco})$, con $\theta_{vco}$ la fase de su propio VCO -- la estimación actual de $\theta$. Repitiendo el cálculo de la Sección \ref{sec:deteccion coherente} (identidades trigonométricas, descartando los términos en $4\pi f_c t$ que el filtro de recepción elimina) pero dejando $\theta_e:=\theta-\theta_{vco}$ explícito en vez de suponerlo nulo:
\begin{equation}\label{eq:xi xq con error de fase}
    x_i(t) \approx m(t)\cos\theta_e, \qquad\qquad x_q(t) \approx m(t)\sin\theta_e, \qquad\text{con } m(t):=\sum_k a_k\, p(t-kT_s).
\end{equation}
Si $\theta_e=0$, recuperamos exactamente lo que la Sección \ref{sec:deteccion coherente} ya usaba ($x_i\approx m(t)$, $x_q\approx 0$). Si $\theta_e\neq 0$, la energía de $m(t)$ se reparte entre $x_i$ y $x_q$: la constelación (en este caso, real, sobre el eje horizontal) aparece \emph{rotada} un ángulo $\theta_e$.

Necesitamos un detector de fase que produzca una señal proporcional a $\theta_e$ a partir de $x_i(t)$ y $x_q(t)$. Ninguno de los dos por separado sirve ($\cos\theta_e$ es par en $\theta_e$, no distingue signo); pero su \emph{producto} sí:
\begin{equation}\label{eq:error costas}
    e(t) := x_i(t)\,x_q(t) \approx m(t)^2\cos\theta_e\sin\theta_e = \frac{m(t)^2}{2}\sin(2\theta_e).
\end{equation}
Para $\theta_e$ chico, $\sin(2\theta_e)\approx 2\theta_e$, y como $m(t)^2\geqslant 0$ siempre, $e(t)$ es -- en promedio sobre varios símbolos, ya que $m(t)^2$ fluctúa símbolo a símbolo pero es siempre no negativo -- proporcional a $\theta_e$ con el signo correcto: exactamente el detector de fase que el PLL genérico de la sección anterior necesita, con $K_d\propto \mathrm E[m(t)^2]$. Este esquema -- multiplicar $x_i$ por $x_q$ para generar el error de fase -- se conoce como \emph{lazo de Costas}, y es la forma estándar de recuperar la portadora para BPSK.

\subsection*{Generalización a M-QAM: decisión realimentada}

Para constelaciones complejas ($M$-QAM, $M$-PSK con $M>2$), conviene trabajar directamente con el símbolo complejo $c_k=a_k+jb_k$ y su versión recibida (rotada) $z_k := x_i(kT_s)+jx_q(kT_s)$. De \eqref{eq:xi xq con error de fase} generalizada a símbolos complejos, $z_k \approx c_k\,\e^{j\theta_e}$: el símbolo recibido es el transmitido, rotado el ángulo de error. El producto que usamos en \eqref{eq:error costas} para BPSK se generaliza reemplazando el símbolo transmitido (que en recepción no conocemos) por la \emph{decisión} del detector, $\hat c_k$ -- la estimación de qué símbolo se transmitió, tomada sobre la constelación completa (Sección \ref{sec:modulacion pasabanda m aria}):
\begin{equation}\label{eq:error decision directed}
    e_k := \mathrm{Im}\left\{z_k\, \hat c_k^*\right\} \approx \mathrm{Im}\left\{|c_k|^2\, \e^{j\theta_e}\right\} = |c_k|^2\sin\theta_e \approx |c_k|^2\,\theta_e,
\end{equation}
donde la aproximación $\hat c_k\approx c_k$ vale si la decisión es mayormente correcta -- lo cual, a su vez, requiere que el lazo ya esté razonablemente cerca del enganche. Este esquema -- llamado \emph{decision-directed}, porque usa las propias decisiones del receptor en vez de un símbolo conocido de antemano -- es el mecanismo que efectivamente se usa en la práctica para $M$-QAM: para $\theta_e=0$, $e_k=0$ (nada que corregir); si el lazo se aparta, $e_k$ crece con el error y lo empuja de vuelta. Nótese la condición de funcionamiento que distingue a este lazo del de Costas: se retroalimenta de sus \emph{propios aciertos}, con el riesgo de que, si arranca demasiado desenganchado, las decisiones sean mayormente erróneas y el lazo diverja en vez de converger. Por eso, en la práctica, el lazo arranca con una secuencia de símbolos \emph{verdaderamente} conocidos (no decididos) antes de pasar a régimen decision-directed -- un motivo más, que se suma a los de la Sección \ref{sec:sincronizacion trama}, para necesitar un preámbulo.

\subsection*{Ambigüedad de fase}

Tanto \eqref{eq:error costas} como \eqref{eq:error decision directed} tienen un problema sutil pero importante: no se anulan solo en $\theta_e=0$. Para el lazo de Costas, $\sin(2\theta_e)=0$ también en $\theta_e=\pi$ -- y es un cero del mismo tipo (mismo signo de la derivada), es decir, un punto de enganche igual de estable que $\theta_e=0$. El lazo puede converger a cualquiera de los dos, según las condiciones iniciales o el ruido durante la adquisición. Si engancha en $\theta_e=\pi$, la constelación de BPSK aparece perfectamente alineada sobre el eje horizontal -- pero \emph{invertida}: cada símbolo $a_k$ se decodifica como $-a_k$. La constelación se ve bien; todos los bits salen mal.

Esto no es un defecto de diseño del lazo, sino una consecuencia de la simetría de la propia constelación: BPSK es invariante ante una rotación de $\pi$ (el punto $-A$ rotado $\pi$ es $A$, y viceversa), así que \emph{ninguna} medida basada solo en la forma de la constelación recibida puede distinguir $\theta_e=0$ de $\theta_e=\pi$. Lo mismo ocurre, de forma más general, para cualquier constelación con simetría rotacional de orden $m$ (invariante ante rotaciones de $2\pi/m$): QPSK y $4$-QAM tienen ambigüedad de orden $4$ (múltiplos de $\pi/2$), y en general una constelación $M$-PSK tiene ambigüedad de orden $M$. La resolución no puede venir del lazo de fase en sí -- viene de afuera: o bien codificando la información \emph{diferencialmente} (cada símbolo codifica un \emph{cambio} respecto del anterior, insensible a una rotación global constante), o bien -- la solución que desarrollamos en la Sección \ref{sec:sincronizacion trama} -- insertando un preámbulo de símbolos conocidos que permite, por comparación directa, detectar y corregir cuál de las rotaciones ambiguas ocurrió.

\subsection*{Temporización de símbolo}

Supongamos ahora resuelto el sincronismo de portadora: la constelación recibida ya no rota. Queda el segundo problema, independiente del primero: el reloj que dispara el muestreador del receptor tampoco coincide con el que generó los símbolos en el transmisor. Llamemos $\tau$ al error de temporización -- el receptor muestrea en $kT_s+\tau$ en vez de en el instante correcto $kT_s$.

La Sección \ref{sec:nyquist no ideal} construyó los pulsos de Nyquist precisamente para que $p(kT_s)=0$ para todo $k\neq 0$: la ISI se cancela \emph{exactamente en el instante de muestreo correcto}. Pero esa cancelación es frágil ante un error de temporización: como muestra la Figura \ref{fig:coseno alzado}, $p(t)$ cae rápido a los costados de $t=0$, así que un $\tau$ chico ya reintroduce ISI de los símbolos vecinos -- el mismo mecanismo, ahora por un desajuste de reloj en vez de un canal no ideal.

Necesitamos, otra vez, un detector de error -- ahora sensible a $\tau$ -- para cerrar un lazo como el de la Sección \ref{sec:sincronizacion} (Figura \ref{fig:pll generico}), pero controlando la fase del reloj de muestreo en vez de la del oscilador de portadora. La idea más directa (\emph{early-late}) es comparar dos muestras tomadas ligeramente antes y después del instante candidato, y usar su diferencia como señal de error: si el muestreo está temprano, una muestra ve más energía del símbolo siguiente; si está tarde, del anterior. El \emph{detector de Gardner} formaliza esto de manera eficiente, con una sola muestra adicional por símbolo: además de la muestra ``a tiempo'' $y_k:=y(kT_s+\tau)$, toma una muestra en el punto medio entre símbolos, $y_{k-1/2}:=y\big((k-\tfrac12)T_s+\tau\big)$, y forma
\begin{equation}\label{eq:error gardner}
    e_k := y_{k-1/2}\,\big(y_k - y_{k-1}\big).
\end{equation}

Veamos por qué \eqref{eq:error gardner} es sensible a $\tau$, aprovechando lo que ya sabemos del pulso $p(t)$: es par (simétrico) y verifica $p(0)=1$, $p(T_s)=0$. Quedémonos con los dos símbolos vecinos dominantes en cada muestra (aproximación de primer orden) y expandamos a primer orden en $\tau$, usando $p(\tau)\approx p(0)=1$ (el término lineal se anula porque $p$ es par) y $p(T_s+\tau)\approx p(T_s)+\tau p'(T_s) = \tau p'(T_s)$:
\begin{equation*}
    y_k \approx a_k,\qquad y_{k-1}\approx a_{k-1}, \qquad y_{k-1/2}\approx (a_{k-1}+a_k)\,p(T_s/2) + \tau\, p'(T_s/2)\,(a_{k-1}-a_k),
\end{equation*}
donde para $y_{k-1/2}$ usamos $p(T_s/2+\tau)\approx p(T_s/2)+\tau p'(T_s/2)$ y, por la paridad de $p$, $p(-T_s/2+\tau)\approx p(T_s/2)-\tau p'(T_s/2)$. Sustituyendo en \eqref{eq:error gardner} y tomando valor esperado sobre símbolos independientes de media nula y varianza $\sigma^2$ (usando $\mathrm E[a_k^2-a_{k-1}^2]=0$ y $\mathrm E[(a_k-a_{k-1})^2]=2\sigma^2$):
\begin{equation}\label{eq:gardner lineal}
    \mathrm E[e_k] \approx -2\sigma^2\, p'(T_s/2)\;\tau.
\end{equation}
Es decir: en promedio, $e_k$ resulta \emph{lineal} en el error de temporización $\tau$, con un signo determinado por $p'(T_s/2)$ -- típicamente distinto de cero para un pulso de Nyquist (el pulso todavía está cayendo hacia su primer cero en $T_s/2$). Esta es exactamente la característica de detector de fase que el lazo de la Sección \ref{sec:sincronizacion} necesita: basta usar $e_k$ como entrada al mismo lazo de primer orden ya derivado (Ecuación \eqref{eq:pll edo}), ahora controlando la fase del reloj de muestreo en vez de la del VCO de portadora, y la convergencia exponencial ya demostrada aplica sin cambios.

\begin{remark}
La Ecuación \eqref{eq:gardner lineal} retiene solo el aporte de los símbolos inmediatamente vecinos a cada muestra; un cálculo completo, incluyendo los símbolos más alejados que también contribuyen (el pulso de Nyquist no se anula \emph{estrictamente} fuera de $kT_s$ salvo en esos puntos), corrige la constante de proporcionalidad pero no cambia la conclusión cualitativa -- linealidad en $\tau$ y el signo del detector -- que es lo que importa para que el lazo enganche.
\end{remark}

\subsection*{Sincronización de trama}\label{sec:sincronizacion trama}

Con portadora y reloj de símbolo ya enganchados, el receptor recibe un flujo continuo y correctamente muestreado de símbolos -- pero sigue sin saber \emph{dónde}, dentro de ese flujo, empieza el mensaje que le interesa. Y, como vimos, puede no saber tampoco cuál de las rotaciones de fase ambiguas enganchó el lazo de portadora. Ambos problemas se resuelven con la misma herramienta: un \emph{preámbulo}, una secuencia de símbolos $s_0,\ldots,s_{L-1}$ \emph{conocida de antemano por transmisor y receptor}, transmitida al comienzo de cada trama.

El receptor busca el preámbulo \emph{correlacionando} la señal recibida contra la secuencia conocida, desplazando una ventana de longitud $L$ sobre el flujo de símbolos:
\begin{equation*}
    R[n] := \sum_{m=0}^{L-1} y_{n+m}\, s_m^*,
\end{equation*}
y evaluando $|R[n]|$ para cada posición candidata $n$. Si el preámbulo efectivamente empieza en $n=n_0$, entonces $y_{n_0+m}\approx s_m$ (a menos de ruido y, eventualmente, la rotación de fase ambigua) para $m=0,\ldots,L-1$, y $R[n_0]\approx\sum_m |s_m|^2$ -- la energía del preámbulo, un valor sustancialmente mayor que $R[n]$ para cualquier otro $n$, donde se están correlacionando símbolos de información (esencialmente aleatorios) contra la secuencia conocida. Buscar el máximo de $|R[n]|$ localiza entonces el inicio de la trama.

Para que este pico sea nítido -- y no se confunda con un máximo espurio en otra posición --, la secuencia $\{s_m\}$ debe diseñarse con \emph{buena autocorrelación}: un pico marcado en el desplazamiento nulo y valores chicos en cualquier otro desplazamiento. Existen familias de secuencias diseñadas explícitamente con este criterio; un ejemplo clásico son los \emph{códigos de Barker}, secuencias binarias ($s_m=\pm 1$) cuya autocorrelación fuera del pico principal está acotada por $1$ en valor absoluto. Por ejemplo, la secuencia de Barker de longitud $7$,
\begin{equation*}
    (s_0,\ldots,s_6) = (+1,+1,+1,-1,-1,+1,-1),
\end{equation*}
tiene autocorrelación $R[0]=7$ en el pico, y $|R[n]|\leqslant 1$ para todo $n\neq 0$ -- una relación pico a lóbulo de $7$ a $1$, excelente para una secuencia tan corta.

El preámbulo cumple, entonces, una doble función, que conecta directamente con lo desarrollado en esta sección:
\begin{itemize}
    \item \textbf{Marca el inicio de la trama}, vía el pico de correlación recién descripto.
    \item \textbf{Resuelve la ambigüedad de fase} de la Sección \ref{sec:sincronizacion}: como el receptor conoce $s_m$, compara la fase del preámbulo \emph{recibido} contra la esperada y determina directamente cuál de las $m$ rotaciones ambiguas enganchó el lazo de Costas (o decision-directed), corrigiéndola antes de decodificar el resto de la trama. De paso, al ser símbolos genuinamente conocidos (no decididos), es exactamente la secuencia de arranque que el lazo decision-directed de la Sección \ref{sec:sincronizacion} necesitaba antes de poder confiar en sus propias decisiones.
\end{itemize}

Con esto se cierra el círculo de este capítulo: la ecualización le devuelve al pulso su forma libre de ISI (Sección \ref{sec:ecualizacion}), y el sincronismo de portadora, de símbolo y de trama garantizan, en conjunto, que el receptor demodula en la fase correcta, muestrea en el instante correcto, y sabe dónde empieza a leer -- las tres suposiciones que, desde el Capítulo \ref{chapter:deteccion ruido}, veníamos dando por sentado.




